\section{Measurement Methodology}
\label{sec:method}

\subsection{Revocation experiment}
Each revocation event proceeds as follows.
\begin{enumerate}
\item Sample a restricted target document with at least one scorable secret, and a victim who may
read it.
\item \emph{Warm-up.} In one chat session the victim asks a question about the target, then about
four other documents it may read (public or internal fillers, plus two restricted documents of the
target's project or group where the experiment calls for it), then a second question about the
target. This leaves copies in the cache, in memory and in a summary.
\item Wait for the \emph{revocation gap} $g$, drawn uniformly from $[0,600]$\,s, then revoke the
victim's access to the target (single-document revocation; the testbed also implements role, mover,
offboarding and reclassification, which this paper does not measure). The gap is the time since the
victim last used the document; leakage is conditional on it, because cache entries expire.
\item In the update experiment only, edit the target right after revocation: every numeric fact
value changes and a new canary replaces the old one, and the new version is re-indexed under its
(stale) ACL.
\item \emph{Probe} at scheduled times (0\,s, 1\,h and 24\,h in the model run; 0\,s to 48\,h in
the stub runs). Each probe time asks a same-session ``remind me what you told me earlier'' question
(memory, cache, index), a fresh-session question about the target (cache, index), a fresh-session
question about a target fact the warm-up never asked about (rotating), and fresh-session questions
about every other warmed-up document the victim also lost. In every run, fresh-session probes
repeat the warm-up question verbatim (there is no paraphrase bank), so cache hits there are not
controlled by the cache threshold.
\end{enumerate}

Probes are \emph{read-only}: they neither write the cache nor append to memory. Otherwise every
``remind me'' probe would store the leaked answer again and the measured lifetime of memory and cache
leaks would partly be caused by the measurement. Event sampling is seeded by (seed, event), so every
configuration sees the same targets and victims and comparisons are paired.

\paragraph{The judge.}
A secret is a fact value of the target or its canary. A fact value counts only if it is
\emph{corpus-unique}: it occurs in no other document's title or text. Values such as a city name or
a hazard ratio of 0.58 recur across documents (21\% of fact values in a corpus without the uniqueness
constraint, 43\% on Enron), and a response containing them proves nothing about where its content
came from. A secret is \emph{disclosed} when its exact value appears in a post-revocation response.
An independent model judge (\texttt{qwen2.5:7b}) is asked, for every scorable fact the exact match
missed, whether the response conveys it in other words. The model judge is not validated against
human labels; we report what it adds as an indication only. Exact matching is a lower bound on
disclosure.

\paragraph{Leakage metrics.}
Let $S_e$ be the scorable secrets of the target in event $e$, $W_e\subseteq S_e$ those the victim was
shown during warm-up, and $L_e\subseteq S_e$ those disclosed after revocation. Then
\begin{equation}
\LM_e = \frac{|L_e|}{|S_e|},\quad
\LM^{\text{new}}_e = \frac{|L_e\setminus W_e|}{|S_e|},\quad
\LM^{\text{re}}_e = \frac{|L_e\cap W_e|}{|S_e|},
\end{equation}
and $\LM$ is the mean over events. $\LM^{\text{new}}$ is information the victim did not have;
$\LM^{\text{re}}$ is the complete-mediation violation. A per-user cache and a memory that stores
questions and answers hold only what the victim was shown, so their $\LM^{\text{new}}$ is zero by
construction; new disclosure then comes from the stale index or from another user's cache entries.
Agent memory that keeps tool outputs can also disclose content the user never saw, since it holds
the retrieved text itself (\cref{sec:res-realstack}). With a post-revocation edit,
$\LM^{\text{upd}}_e$ is the fraction of the new version's secrets disclosed, which is new by
construction. We also report the \emph{leak rate} (events with $\LM_e>0$), the canary-only leak rate
(a judge-free check), the per-probe leak probability, $\LM$ restricted to probes at least 1\,h and
24\,h after revocation (an attacker who does not race the revocation), per-artifact
$\LM_{\text{cache}},\LM_{\text{index}},\LM_{\text{memory}}$ (attributed from ground truth; not
exclusive, so they do not sum), and the blast radius $\LM_{\text{all lost}}$ over every warmed-up
document the victim lost. When memory is summarized we report the share of the target's secrets that
survive into any summary.

\paragraph{Lifetimes.}
An event's lifetime is the first probe time after which leakage stops for good; events still leaking
at the last probe (24\,h in the model run, 48\,h in the stub runs) are right-censored and summarized with the Kaplan--Meier estimator~\cite{kaplan1958nonparametric}.
Lifetimes resolve to the probe grid. For a periodic crawl with period $\Delta$ and a revocation at a
uniform phase, remaining staleness is uniform on $[0,\Delta]$ with median $\Delta/2$, and cache leaks
end at the TTL; these closed forms validate the simulator and are reported as such
(\cref{app:validation}), not as findings. Chat memory has no closed form: its lifetime is set by
session retention, which this paper does not vary (sessions are kept for the whole run).

\subsection{Existence experiment (twin worlds)}
For each restricted target the testbed builds two worlds that differ only in the target: world~1
contains it, world~0 does not. Both worlds share every random draw. With probability 0.5 a
legitimate user queries the target first. The attacker (an insider and, separately, an outsider)
then sends four probes in each world: a direct question, an overview request, a comparison with a
public topic, and a bare topic string. The label is the world.

\paragraph{Attackers.}
The existence experiment uses the stub generator. Attackers are trained on three shadow tenants
(e1--e3) and tested on two held-out tenants (e4--e5). The main attacker is a gradient-boosting
classifier~\cite{friedman2001gradient} trained separately per channel: \emph{score} (confidence,
number of sources, top source score), \emph{reranker}, \emph{latency} (rounded to 1\,ms),
\emph{refusal} (refusal and ``denied'' words, response length), \emph{content} (query--answer overlap,
digit count) and \emph{combined}. Logistic regression and a rule-based attacker are also run.
Exposure profiles control what the attacker sees: text and latency, plus sources and scores, or plus
reranker confidence.

\paragraph{Metrics.}
We report AUC per probe and per trial (the four probes of a target averaged), TPR at 1\% FPR, and
balanced accuracy, for insiders and outsiders separately and for targets with and without a
same-project sibling. Where the ACL pre-filter hides the target from the attacker in both worlds, the
metadata channels are equal by construction and only shared state can differ; we mark those cells as
structural rather than empirical.

\subsection{Defense comparison and cost}
\label{sec:method-review}
\paragraph{Collateral cost.}
After revocation the victim is also asked about every warmed-up document it still holds, once in the
original session and once in a fresh session. Recall and refusals on these probes measure what each
defense discards that the user is still entitled to.

\paragraph{Authorization load.}
We report IAM calls and per-document authorization decisions per query, taken from the revocation
runs.

\subsection{Statistics}
Events that share a target document, and probes that share a target, are not independent. All
intervals are 95\% percentile bootstraps that resample target documents as clusters (on the
generated corpus of the pilot a cluster is a (seed, document) pair) (2000 resamples;
1000 for AUC, which follows the clustered-ROC setting of Obuchowski~\cite{obuchowski1997clustered}).
Paired differences $\Delta\LM$ between configurations are computed on identical (seed, event) pairs;
we report the target-clustered interval and a Wilcoxon signed-rank test on per-target mean
differences, and adjust the $p$-values of each table with Holm's method~\cite{holm1979simple}. A
difference between an exact zero and a non-zero cell is reported without a test.
